\section*{Anexo 4-E — Trazabilidad funcional}
\phantomsection\label{fuente-06-anexo-4-1-e-trazabilidad-funcional}
\phantomsection
\label{anx:E}
\addcontentsline{toc}{section}{Anexo 4-E — Trazabilidad funcional}
La tabla sigue cada grupo de requisitos funcionales hasta el módulo responsable, las capas que lo realizan y la evidencia que permite verificarlo.

\begin{tablalafrox}{Trazabilidad funcional de la arquitectura lógica}{tab:traza-logica}{F{0.16} F{0.26} F{0.26} Y}{\cab{Requisito} & \cab{Responsable} & \cab{Capas implicadas} & \cab{Evidencia verificable}}
RF-01 & M1 Recepción & 1, 4, 5 y 6 & Lote GS1 recibido y publicado. \\
RF-02 & M2 Inventario & 4, 5 y 6 & Reserva y FEFO por sitio. \\
RF-03 & M3 Preventa & 1, 3, 4 y 6 & Pedido único tras reconexión. \\
RF-04 & M4 Rutas & 1, 4 y 5 & Ruta revisada en menos de 20 min. \\
RF-05 & M5 Preparación & 1, 4, 5 y 6 & Lectura y carga por lote. \\
RF-06 & M6 Reparto & 1, 4, 5 y 6 & POD unido a entrega. \\
RF-07 & M7 Cobranza y rendición & 1, 4 y 5 & Descuadre y aprobación segregados. \\
RF-08 & M8 Devoluciones y envases & 1, 4, 5 y 6 & Saldo y causal conciliados. \\
 RF-09  &  M9 Calidad y trazabilidad  &  1, 2, 3, 4, 5, 6 y 7  &  Bloqueo local, evaluación por Calidad con permiso vigente y decisión durable a M2/M5.  \\
RF-10 & M2 Inventario & 4, 5 y 6 & Sugerencia por SKU y proveedor con promociones comprometidas, aprobada por abastecimiento y traspasada al ERP por la ACL de INT-06. \\
RF-11 & M10 Analítica & 4, 5 y 6 & OTIF y costo trazables al evento. \\
RF-12 & M11 Canal moderno & 3, 4 y 5 & Pedido EDI normalizado. \\
RF-13 & Base compartida & Transversal & Entregables de arquitectura y contratos trazables a los requisitos. \\
 RF-14  &  M12 Telemetría  &  2, 4 y 5  &  Ruta real correlacionada.  \\

 RF-15  &  Identidad y acceso  &  1, 3 y 7  &  Autoridad Keycloak, sesiones, permisos y verificador local, 4.1.3.7 y Anexo 4-N.  \\

 RF-16  &  Observabilidad  &  4, 5 y 8  &  Correlación, métricas, trazas y buffer local, 4.1.3.8 y Anexo 4-N.  \\

 RF-17  &  Administración e interacción transversal  &  1, 3, 4, 6 y 7  &  Interfaces autorizadas, parametrización versionada y funciones comunes, 4.1.3.1 y 4.1.3.4.  \\

 RF-18  &  Firma, documentos y notificaciones  &  1, 4, 5, 6 y 7  &  Firma por acto, POD, DTE, acuse y entrega de avisos; 4.1.17, INT-11 y Anexo 4-U.  \\

 RF-19  &  Gobierno y clasificación de información  &  4, 6 y 7  &  Clasificación, permisos, cifrado y retención, 4.1.3.6–4.1.3.7 y Anexo 4-R.  \\
\end{tablalafrox}

Fuente: elaboración a partir del catálogo de requerimientos del capítulo 3 y de la Tabla~\ref{tab:modulos-funcionales}. Diseño propuesto: (Subdocumento 4, anexos, diseño de este apartado, pp.~\pageref*{fuente-06-anexo-4-1-e-trazabilidad-funcional}--\pageref*{fuente-06-anexo-4-1-e-trazabilidad-funcional-fin}).

La tabla agrupa familias para facilitar la lectura. La trazabilidad individual de RF, RNF y RT se conserva en el Formulario T-12, que identifica el componente y la sección que desarrolla cada requisito. Los anexos 4-D, 4-F y 4-N aportan responsabilidad, contrato, capas y realización física. T-11 identifica el recurso ofertado. El estado de cumplimiento de T-12 no demuestra por sí solo una prueba ejecutada.

Los RNF se siguen por la capacidad que condicionan: desempeño y concurrencia en 4.2.6 y Anexo 4-T. Autonomía y sincronización en 4.1.15–4.1.16 y anexos 4-J/4-K, identidad, seguridad y auditoría en 4.1.3.7 y anexos 4-Q/4-R, disponibilidad y recuperación en 4.2.4–4.3.2. Observabilidad en 4.1.3.8. Para soporte y capacitación, SD4 declara los recursos y dependencias arquitectónicas, mientras sus planes y evidencias se desarrollan en los subdocumentos respectivos. Así se conserva la correspondencia sin atribuir al capítulo 4 responsabilidades de ejecución ajenas.

\label{fuente-06-anexo-4-1-e-trazabilidad-funcional-fin}
